\subsection*{Variables and Types}

Croqtile introduces four fundamental type categories: scalar-type, spanned-type, integer-tuple-type (ituple-type) and bounded-type. These types serve for different purposes.

\begin{itemize}[leftmargin=1.6em]
\item \textbf{Scalar Type}. It is designed to fulfill the requirement of program control. It consists of \emph{Integer Type} and \emph{Boolean Type}.
\item \textbf{Spanned Type}. It represents the data type (normally the tensor) for computation. Apart from referencing the raw data, a \emph{spanned type} also associates data with a shape represented by multi-dimensional ranges. The 'multi-dimensional ranges' is named 'mdspan'. It is useful for tiling, etc, which is introduced soon.
\item \textbf{Integer Tuple (I-Tuple) Type}. It represents a group of integer values. A common usage of \emph{i-tuple} is to introduce the bound of multi-dimensional data.
\item \textbf{Bounded (Integer/ITuple) Type}. This is the special type that is used to simplify data (sub-zone) reference. And it works with loop construct to iteratively process data.
\end{itemize}

Scalar and \texttt{ituple} values resemble familiar constructs in general-purpose languages. Spanned and bounded types are specific to Croqtile's domain; the following sections explain their roles.
